\itemhead{习题4.2 · 第8题（教材p.114）}{XJ-B1-C04-HW-P114-Y0667}
\label{XJ-B1-C04-HW-P114-Y0667}
{\footnotesize 湘教版教材p.114 / PDF第121页 / 习题4.2 / 第8题}
\field{原题}{已知 $0<a<1$，$b<-1$，函数 $y=a^x+b$ 的图象不经过哪个象限？}
\field{参考答案}{不经过第一象限。$x\ge0$ 时 $a^x\le1$，故 $a^x+b<0$；在负半轴既有正值也有负值，因此其余三个象限均经过。}
{\small \field{关联知识}{E01 指数函数的概念、图象与性质（教材帮p.150、151 / PDF159、160）；F09 图象平移、伸缩与对称（教材帮p.86、94 / PDF95、103）}
\field{方法与易错点}{先查底数条件，利用已知点定底数，再据单调性及图象求值或范围。 易错：只要求底数非负；漏排1；将 $a^x+b$ 当标准指数函数且仍用原值域。}
\field{答案依据}{答案册相关题组：PDF第32页（册内页31）；按当前题库原题逐项重算，不按旧版题号机械对套。}}
